% !PRINTED 355
\weberpage{355}

One can now apply the method of interpolation to this in order to obtain a more accurate value of $A$, by distributing the total error of $0{,}236$ over the two values of $A$ in proportion to the partial errors, that is, by putting
\[
A = 0{,}2 + \frac{0{,}1}{0{,}236}\,0{,}113 = 0{,}247\ldots
\]

With the help of a seven-place table of \gesperrt{Zach} one now obtains:
\[
\begin{array}{c|c|c|c}
A & B & 2A - B & \text{Error} \\
\hline
0{,}247 & 0{,}1948581 & 0{,}2991419 & -0{,}0018881 \\
0{,}248 & 0{,}1944969 & 0{,}3015031 & +0{,}0004731
\end{array}
\]
and by a further application of the interpolation
\[
A = 0{,}2477996
\]
or % [sic: the page prints "odre" for "oder"]
\[
x = 1{,}769292.
\]

\begin{center}
\S.\,114.

Computation of the imaginary roots of a trinomial equation.
\end{center}

\gesperrt{Gauss} has also given a procedure for the computation of the imaginary roots of a trinomial equation, which we shall briefly discuss here. We make the assumption of real coefficients, although the method is applicable to the general case as well. We shall put the equation in the form:
\[
(1)\quad x^{m+n} - ex^m - g = 0,
\]
but here we need not restrict ourselves to positive values of $e$ and $g$.

We set
\[
x = re^{i\vartheta}
\]
and, by setting the imaginary part in (1) equal to zero by itself, we obtain
\[
(2)\quad r^n = \frac{e\sin m\vartheta}{\sin(n+m)\vartheta};
\]
%CONT>>
